\chapter{Introducción}

El presente proyecto de grado se centró en abordar los desafíos que enfrentan los talleres industriales en la gestión eficiente de sus operaciones, específicamente en la asignación de tareas dentro de un sistema de Planificación de Recursos Empresariales (ERP). La asignación tradicional de tareas, frecuentemente realizada de forma manual, puede conducir a ineficiencias significativas, tales como la sobrecarga de operarios, tiempos muertos en la producción y retrasos en las entregas, afectando directamente la productividad y la rentabilidad del taller.

En este contexto, la automatización de la asignación de tareas emerge como una solución crucial para optimizar el uso de los recursos y mejorar el flujo de trabajo. Este proyecto propuso el desarrollo de un modelo de automatización basado en algoritmos genéticos, una técnica de inteligencia artificial que ha demostrado su eficacia en la resolución de problemas complejos de optimización. Al implementar algoritmos genéticos en el módulo de gestión de órdenes de trabajo de un ERP, se buscó automatizar la distribución de tareas considerando variables clave como las capacidades técnicas de los operarios, la disponibilidad de insumos y la carga de trabajo actual.

Para el desarrollo del proyecto, se llevó a cabo un proceso metodológico que incluyó el análisis de los procesos actuales de asignación de tareas en Rectificadora Recti-Ram, la recopilación y estructuración de datos operacionales, el diseño e implementación de un algoritmo genético adaptado a las necesidades específicas del taller, y la evaluación comparativa del modelo propuesto frente a la heurística clásica \textit{Shortest Processing Time} (SPT), ampliamente utilizada en el campo de la ingeniería industrial. Los resultados obtenidos demostraron que el algoritmo genético desarrollado logró una reducción del 15.6\% en el makespan promedio (de 186.53 a 157.37 unidades de tiempo) y una disminución del 14.3\% en la desviación estándar de carga entre operarios (de 35.11 a 30.10), logrando cronogramas más compactos, mejor balance de trabajo y mayor estabilidad operativa. Estas mejoras validaron la eficacia y aplicabilidad del enfoque propuesto en un entorno industrial real, demostrando que el algoritmo genético no solo es competitivo frente a métodos tradicionales, sino que los supera en los aspectos más relevantes para la operación diaria del taller.

La investigación se llevó a cabo en colaboración con la empresa Rectificadora Recti-Ram, lo que proporcionó un contexto real y datos relevantes para el desarrollo y la validación del modelo.

A través de este proyecto, se contribuyó al avance en la automatización de los talleres industriales, proporcionando una herramienta que permite mejorar la eficiencia operativa, reducir los costos y aumentar la satisfacción del cliente mediante la optimización de los procesos de asignación de tareas.